\begin{formulaBox}{Binomialverteilung}{ML I}
Bei $n$ unabhängigen Bernoulli-Versuchen mit derselben
Erfolgswahrscheinlichkeit $\pi$ zählt $X$ die Erfolge:
$X\sim B(n,\pi)$ mit $0\leq\pi\leq1$.

\begin{multicols}{2}
\subsection*{Wahrscheinlichkeitsmasse}
Für $k\in\{0,1,\dots,n\}$ gilt:
\[
P(X=k)=\binom{n}{k}\pi^k(1-\pi)^{n-k}
\]

\subsection*{Binomialkoeffizient}
\[
\binom{n}{k}=\frac{n!}{k!(n-k)!},\qquad 0!=1
\]

\columnbreak
\subsection*{Erwartungswert und Varianz}
\begin{align*}
\operatorname{E}(X)&=n\pi\\
\operatorname{VAR}(X)&=n\pi(1-\pi)
\end{align*}
Der Stichprobenanteil ist $\hat{\pi}=X/n$.
Die zugehörigen Standardfehler stehen im folgenden Abschnitt.
\end{multicols}
\end{formulaBox}
